\noindent WireGuard is a widely adopted, modern, free and open-source VPN protocol known for its high performance, minimal codebase, and strong security. Like all VPN protocols relying on classical public-key cryptography, it faces a growing quantum threat in which adversaries are already collecting encrypted traffic today, intending to decrypt it once sufficiently powerful quantum computers become available, a strategy known as ``harvest now, decrypt later''. This thesis evaluates the integration of two NIST-selected post-quantum Key Encapsulation Mechanisms (KEMs): ML-KEM and Hamming Quasi-Cyclic (HQC) into WireGuard's handshake protocol. Both hybrid and pure-KEM handshake variants were designed and implemented in wireguard-go, the official userspace Go implementation of WireGuard, and benchmarked across handshake latency, message size, and per-peer storage requirements. Results show that ML-KEM integrates well as the hybrid ML-KEM-768 variant fits within a single IPv4 packet, adds modest overhead to handshake latency, and keeps storage requirements manageable at scale. Pure ML-KEM variants achieve full post-quantum authentication at the cost of larger messages requiring fragmentation. HQC, by contrast, produces messages large enough to force heavy fragmentation per handshake, resulting in latencies that render it incompatible with WireGuard's design in its current form. The findings suggest that a hybrid ML-KEM-768 deployment is achievable within existing WireGuard infrastructure today, offering protection against harvest now, decrypt later attacks without replacing the current static key infrastructure or causing significant performance degradation. \\

\noindent \textbf{Keywords:} WireGuard, VPN, Post-Quantum, PQ, Post-Quantum Cryptography, PQC, ML-KEM, HQC, Key Encapsulation Mechanism, KEM, Hybrid Cryptography